\section{Domain interfaces and reusable proof recipes}
\label{sec:domains}

The framework is deliberately domain-neutral, but its value is clearest when the
abstract nodes are translated into concrete mathematical recipes.  This section gives
such recipes and identifies the remaining input for each area.

\subsection{Finite algebra and combinatorics}

Exact finite algebra is the easiest case.  Encode tuples, lists, matrices, polynomials,
permutations, graphs, and finite sets by natural numbers.  Operations are primitive
recursive, and identities are equations or bounded formulas.  The standard recipe is:

\begin{enumerate}[label=\arabic*.,leftmargin=2.2em]
\item choose a canonical code or a proof-relevant equality certificate;
\item define operations by fixed primitive recursion;
\item prove defining equations and congruence in \(\Prf\,\PRAAx\);
\item Booleanize bounded side conditions;
\item prove the main invariant by QF induction on size, degree, or a loop clock; and
\item use the proof schedule to compose algebraic subroutines.
\end{enumerate}

Determinants by Laplace expansion, adjugates, bounded sums/products, Bell recurrences,
finite group actions, and combinatorial bijections fit directly.  Symbolic polynomial
identities may be proved in a host ring and reflected through a checked object-level
ring lemma, as the repository's \code{ipra\_ring} tactic already demonstrates.

\subsection{Exact coordinate geometry and general relativity identities}

Incidence and metric identities over rational or integer coordinates reduce to
polynomial equalities and guarded divisions.  Degenerate configurations are handled by
explicit nonvanishing hypotheses or by clearing denominators.  For symbolic spacetime
metrics, one fixes a coordinate chart and rational/algebraic parameter coding, computes
inverse matrices, Christoffel symbols, curvature, and field tensors by finite sums, and
proves the resulting polynomial/rational identities under denominator guards.

The framework proves local coordinate identities.  It does not infer global causal or
manifold statements unless chart coverage, transition maps, and topological
certificates are supplied as separate nodes.

\subsection{Algebraic numbers, roots, and spectral data}

A robust root interface returns a finite cloud of rational disks or boxes with:

\begin{itemize}
\item residual bounds;
\item a coverage/count certificate;
\item overlap or continuation edges across precision;
\item symmetric-product or coefficient reconstruction estimates; and
\item separation data only when exact individual labels are requested.
\end{itemize}

Multiple roots and eigenvalue collisions are treated relationally.  Stable outputs
include the unordered multiset, elementary symmetric functions, cluster sums,
characteristic polynomials, and spectral projectors onto isolated clusters.  An
ordered eigenvector or individual branch requires a gap and normalization guard.

Variable degree is handled by one clocked recursion on degree after a fixed coefficient
evaluator has been statically specialized.  Backward dual-horizon planning determines
all coefficient rows needed for repeated deflation and residual verification.

\subsection{Real analysis and calculus}

A function presentation should contain the quantitative data familiar from constructive analysis \cite{bishopbridges1985,troelstravanDalen1988}:

\begin{itemize}
\item a rational sampler or regular-name evaluator;
\item a domain certificate;
\item a continuity, derivative, variation, or holomorphy modulus appropriate to the
theorem;
\item magnitude bounds needed for products and remainders; and
\item finite algorithms for partitions, sums, or polynomial approximants.
\end{itemize}

The fundamental theorem of calculus, Taylor bounds, Euler--Maclaurin, and special-
function identities are then stated as rational inequalities at requested precision.
One computes a partition size or truncation order from the supplied moduli, verifies a
finite sum identity, and bounds the remainder.  No induction is performed on the
predicate ``this is a Cauchy name''; the final precision is frozen and only the finite
rows selected by the justification horizon occur in the invariant.

Qualitative differentiability or holomorphy is not silently converted to a uniform
modulus.  For specially constructed functions, a modulus may be derived from explicit
convexity, coefficient, or derivative bounds and then becomes another compiled node.

\subsection{Complex analysis and contour methods}

Contours are finite polygonal or rational parametrized paths with derivative and
clearance bounds.  Analytic functions are represented on a certified neighborhood by
samplers and uniform error/modulus data.  Cauchy and residue arguments become finite
quadrature plus explicit deformation or local-series estimates.  Poles are guarded by
separation from the contour.  Winding numbers can be returned as exact integers when a
positive distance from zero ensures stability under approximation.

An abstract identity theorem or unrestricted contour deformation is not needed for a
fixed application: backward demand propagation expands the finite tree of local
comparisons and approximation estimates needed by the final contour calculation.

\subsection{Analytic number theory}

Dirichlet polynomials, finite Euler products, summatory functions, and truncated
integrals are finite arithmetic.  The difficult inputs are explicit analytic estimates.
An \(\FRNF\) theorem should state them as contracts with all dependencies visible:
modulus, conductor, field discriminant, strip, height, or zero-free-region parameters.
The finite compiler then:

\begin{enumerate}[label=(\roman*),leftmargin=2.4em]
\item selects truncation and contour parameters;
\item allocates precision backward through the finite expression graph;
\item verifies rational/complex approximations and denominator guards;
\item carries error numerators by QF induction; and
\item emits the final explicit inequality.
\end{enumerate}

This architecture can strengthen reliability and reuse without pretending to improve
the underlying zero-free region.

\subsection{Diophantine geometry and complete enumeration}

Given a family \(X_p\) and a primitive-recursive height bound \(H(p)\), enumerate every
rational/integral code of height at most \(H(p)\), test the defining equations, and
return the filtered list.  The object proof establishes soundness by direct checking
and completeness by the height theorem.  Descent and local solubility certificates may
reduce the search or prove the bound in special families.

Without \(H\), a finite verifier remains useful but the complete-list theorem is
conditional.  A cardinality theorem alone cannot fill the gap by
\cref{barrier:cardinality-cutoff}.

\subsection{Probability}

Finite rational probability laws are exact tables.  Expectations, moments,
convolutions, characteristic functions at rational points, and finite-dimensional
distribution comparisons are bounded sums.  Quantitative limit theorems require
explicit rates such as Berry--Esseen constants, truncation/tail bounds, Lindeberg
moduli, or smoothing inequalities.

For an infinite or real-valued law, the interface contains a primitive-recursive
quantization at precision \(k\), a tail/error bound connecting it to the original law,
and the finite-law theorem.  The proof schedule composes these.  It does not quantify
inside first-order arithmetic over arbitrary measures as higher-type objects.

\subsection{Topology}

A finite simplicial complex is exact combinatorial data, so chains, boundaries, cycles,
homology certificates, and simplicial maps are finite linear algebra.  A represented
space is observed through a tower of finite complexes.  A map returns a finite carrier
cloud rather than a forced vertex at ties.  If every selection from the carrier is
simplicial and any two selections are contiguous, an explicit prism operator gives a
chain homotopy and makes the induced homology map independent of choices.

Jordan-type and dimension arguments then reduce to finite grid parity, Sperner labels,
or chain calculations plus supplied sampling and inverse-modulus data.  The observation
tower unifies the presentation; it does not generate missing continuity moduli.

\subsection{Category theory and algebraic geometry}

For fixed finite presentations, categorical constructions are finite data:
objects/morphisms have codes, composition is a table or fixed program, and diagrams are
bounded lists.  Universal properties are represented by explicit witness maps and
uniqueness derivations.  Yoneda components, finite limits/colimits, adjunction units,
and finite Kan extensions can therefore be expressed in \(\iPRA\) for specified
families.

Quotients should be certified setoids, not opaque host quotients.  Equality in a
quotient is witnessed by a finite certificate---for example an ideal combination,
localization killing exponent, or rewrite normalization.  Affine gluing uses finite
covers, overlap maps, and cocycle certificates.  Scheme-theoretic assertions requiring
unbounded families or non-effective existence remain outside the finite presentation
unless additional producers are supplied.

\subsection{A common audit question}

Across all areas, the decisive question is:

\begin{quote}
For the final numerical challenge \(d\), can every object, query, branch, and proof
obligation actually used by the argument be placed in one primitive-recursively bounded
finite graph, with local correctness derivations and explicit external producers?
\end{quote}

If yes, the proof is a candidate for \(\FRNF\).  If no, the failure identifies either a
missing quantitative theorem or an inappropriate representation.
